% ============================================================================
%  C/12-mach-van — file chương
%  Ngày tạo: 2026-09-07 | Sửa gần nhất: 2026-09-07 | Ôn gần nhất: chưa
%  Dependency: C/11, B/06. Mức độ: cốt lõi.
% ============================================================================
\documentclass[../../english.tex]{subfiles}
\begin{document}
\chapter{Mạch văn: cũ trước, mới sau (Flow and Emphasis)}
\ptrtarget{C/12}\ptrtarget{C/12-mach-van}
\dependency{\ptr{C/11} vì mạch chỉ có nghĩa khi từng câu đã rõ; \ptr{B/06} cho các từ nối và bộ khung lập luận.}

\begin{tomtat}
Chương trước lo từng câu; chương này lo chỗ nối giữa các câu. Nguyên tắc trung tâm cũng chỉ có một, và nó là một trong những ý hữu ích nhất của cả quyển\cite{gopenswan1990}: mỗi vị trí trong câu có một chức năng cố định trong đầu người đọc --- \emph{đầu câu} là chỗ họ tìm mối nối với cái vừa đọc, \emph{cuối câu} là chỗ họ đặt trọng âm và nhớ lâu nhất. Văn khó theo dõi thường là văn đặt thông tin ngược lại: mới ở đầu, cũ ở cuối. Nội dung gồm: quy tắc cũ trước mới sau; vị trí nhấn; hai kiểu chuỗi chủ đề; cách sửa một đoạn rời rạc; và vì sao lời khuyên ``dùng nhiều từ nối'' chữa sai bệnh.
\end{tomtat}

\begin{nentang}
\kn{thông tin cũ}{cái người đọc đã biết --- vừa nói ở câu trước, hoặc là kiến thức chung của họ.}
\kn{thông tin mới}{cái bạn đang muốn thêm vào đầu người đọc trong câu này.}
\kn{vị trí chủ đề}{đầu câu, trước động từ chính --- nơi người đọc tìm ``câu này nói về cái gì và nối với cái gì''.}
\kn{vị trí nhấn}{cuối câu, sau động từ --- nơi người đọc đặt trọng âm và nhớ lâu nhất.}
\kn{chuỗi chủ đề}{cách các chủ ngữ nối tiếp nhau qua một đoạn; hai kiểu chính là lặp lại và móc xích.}
\kn{móc xích}{phần mới ở cuối câu này trở thành phần cũ ở đầu câu sau.}
\end{nentang}

\begin{khoigoi}
\h{Hai đoạn có cùng các câu, cùng các từ nối, mà một đoạn dễ theo và một đoạn rời rạc. Khác biệt nằm ở đâu?}
\h{Vì sao đặt thông tin quan trọng nhất ở đầu câu lại thường là một sai lầm?}
\h{Lời khuyên ``thêm \emph{however}, \emph{moreover}, \emph{therefore}'' để văn mạch lạc sai ở chỗ nào?}
\h{Khi nào nên dùng bị động vì lý do mạch văn, dù câu chủ động ngắn hơn?}
\h{Bạn có một đoạn mà người đọc kêu khó theo. Quy trình sửa gồm những bước gì?}
\end{khoigoi}

Hãy bắt đầu bằng một thí nghiệm nhỏ. Đọc hai đoạn dưới đây; chúng chứa đúng cùng một lượng thông tin.

\begin{capdoi}[Cùng thông tin, hai cách sắp]
\truoc{Motion blur degrades feature extraction. Fewer inliers are produced by the matching stage in such frames. A loss of tracking may follow. Aggressive rotation is the usual cause of blur in handheld sequences.}
\sau{Motion blur degrades feature extraction. Poor features give the matching stage fewer inliers, and too few inliers can cause a loss of tracking. Blur itself usually comes from aggressive rotation in handheld sequences.}
\vico{Bản trên: mỗi câu mở đầu bằng một thông tin mới, nên người đọc không thấy chúng nối với nhau. Bản dưới: cái mới ở cuối câu trước quay lại làm cái cũ ở đầu câu sau.}
\end{capdoi}

Không câu nào trong bản trên sai ngữ pháp, và không thiếu từ nối nào. Cái thiếu là \emph{trật tự}: người đọc bước vào mỗi câu mới mà không có điểm tựa. Đó là câu trả lời cho câu hỏi mở đầu thứ nhất, và là nội dung của cả chương.

\section{Hai vị trí, hai chức năng}

Quan sát nền tảng --- Gopen và Swan trình bày nó rõ nhất\cite{gopenswan1990} --- là: người đọc tiếng Anh không xử lý câu như một túi từ mà gán \emph{chức năng} cho vị trí.

\emph{Đầu câu là vị trí chủ đề.} Khi bắt đầu một câu, người đọc mặc định rằng những từ đầu tiên cho biết câu này nói \emph{về} cái gì, và quan trọng hơn, nó nối với cái gì trong những gì họ vừa đọc. Nếu chỗ đó là một khái niệm hoàn toàn mới, họ mất điểm tựa.

\emph{Cuối câu là vị trí nhấn.} Sau động từ, giọng đọc trong đầu tự nhiên lên trọng âm, và thông tin ở đó được nhớ lâu nhất. Đây là chỗ dành cho cái mới, cho con số bạn muốn người ta nhớ, cho kết luận.

Từ hai quan sát này ra hai quy tắc, và chúng là toàn bộ phần thực hành:

\begin{cannho}
\textbf{Quy tắc một --- cũ trước, mới sau.} Bắt đầu câu bằng thông tin người đọc đã có; kết thúc bằng thông tin bạn đang thêm vào.\\[0.25em]
\textbf{Quy tắc hai --- đặt cái đáng nhớ ở cuối.} Nếu một chi tiết đáng để người đọc nhớ, đừng chôn nó giữa câu.
\end{cannho}

Quy tắc hai trả lời câu hỏi mở đầu thứ hai, và nó đi ngược trực giác của nhiều người --- kể cả lời khuyên ``đặt ý chính lên đầu'' mà bạn có thể đã nghe. Lời khuyên đó đúng ở cấp \emph{đoạn} và cấp \emph{bài} (\ptr{C/13} sẽ nói), nhưng sai ở cấp câu. Trong một câu, đầu câu là chỗ để \emph{bắt tay} với cái đã biết, không phải chỗ để tung tin quan trọng nhất.

\begin{figure}[H]\centering
\begin{tikzpicture}[
  box/.style={draw=steel,fill=bluebg,rounded corners=2pt,inner sep=4pt,font=\footnotesize,align=center,minimum height=0.75cm},
  boxn/.style={box,draw=greendark,fill=greenbg},
  ar/.style={-{Latex[length=2.0mm]},draw=accent,thick},
  lb/.style={font=\scriptsize\itshape,color=tiercol,align=center}]
\node[box,text width=3.5cm] (t1) {\textbf{cũ}\\ Motion blur};
\node[boxn,right=0.12cm of t1,text width=4.6cm] (n1) {\textbf{mới}\\ degrades feature extraction};
\node[box,below=0.85cm of t1.south west,anchor=north west,text width=3.5cm] (t2) {\textbf{cũ}\\ Poor features};
\node[boxn,right=0.12cm of t2,text width=4.6cm] (n2) {\textbf{mới}\\ give fewer inliers};
\node[box,below=0.85cm of t2.south west,anchor=north west,text width=3.5cm] (t3) {\textbf{cũ}\\ Too few inliers};
\node[boxn,right=0.12cm of t3,text width=4.6cm] (n3) {\textbf{mới}\\ can cause loss of tracking};
\draw[ar] (n1.south) .. controls +(0,-0.35) and +(0,0.35) .. (t2.north);
\draw[ar] (n2.south) .. controls +(0,-0.35) and +(0,0.35) .. (t3.north);
\node[lb,right=0.35cm of n2] {móc xích:\\ mới \(\rightarrow\) cũ};
\node[lb,above=0.12cm of t1.north west,anchor=south west] {vị trí chủ đề};
\node[lb,above=0.12cm of n1.north east,anchor=south east] {vị trí nhấn};
\end{tikzpicture}
\caption{Chuỗi móc xích. Phần mới ở cuối câu trước quay lại làm phần cũ ở đầu câu sau, nên người đọc luôn có điểm tựa khi bước vào câu mới.}
\label{fig:mach-van-chuoi}
\end{figure}

\section{Hai kiểu chuỗi chủ đề}

Có hai cách tổ chức các đầu câu qua một đoạn, và chọn đúng kiểu là một quyết định thật chứ không phải thẩm mỹ.

\emph{Kiểu lặp lại: cùng một chủ đề chạy suốt đoạn.} Mọi câu đều bắt đầu bằng cùng một thứ (hoặc cách gọi khác của nó). Dùng khi đoạn mô tả \emph{nhiều mặt của một đối tượng}: \emph{The tracker runs at 30 Hz. It keeps a window of ten keyframes. It discards\ldots}

\emph{Kiểu móc xích: cái mới thành cái cũ.} Mỗi câu bắt đầu bằng phần vừa kết thúc câu trước. Dùng khi đoạn kể một \emph{chuỗi nhân quả hoặc quy trình} --- đúng như ví dụ ở đầu chương.

Đoạn rời rạc thường là đoạn không theo kiểu nào cả: đầu câu nhảy lung tung. Và một lỗi tinh vi hơn là trộn hai kiểu giữa chừng mà không có lý do --- người đọc đã quen với một nhịp rồi bị hẫng.

\begin{mauau}[Mẫu câu nối mạch]
\mc{Móc xích tường minh}{This \textup{[danh từ gói lại ý trước]} explains why \ldots}{\emph{This behaviour}, \emph{This gap} --- gói ý câu trước thành cái cũ}
\mc{Giữ chủ đề}{The tracker \ldots It \ldots Its \ldots}{Kiểu lặp lại; tránh đổi cách gọi giữa chừng}
\mc{Đảo để nối}{Such frames are then discarded.}{Bị động để đưa \emph{such frames} (cũ) lên đầu}
\mc{Đưa cái nhấn về cuối}{The error drops to 2 cm only when the map is scaled correctly.}{Điều kiện quan trọng đặt ở vị trí nhấn}
\mc{Mở đoạn mới}{We now turn to \ldots}{Báo hiệu đứt mạch có chủ ý, khác với đứt mạch do lỗi}
\end{mauau}

\begin{cambay}
Đừng gói ý trước bằng một đại từ trần: \emph{This shows that\ldots} Khi \emph{this} đứng một mình, người đọc phải quay lại đoán nó trỏ vào cái gì --- một ``việc dở dang'' đúng nghĩa ở \ptr{C/11}. Luôn cho \emph{this} một danh từ: \emph{This drop}, \emph{This inconsistency}, \emph{This choice}. Đây là một trong những sửa nhỏ cho hiệu quả lớn nhất, và cũng là chỗ danh từ hoá trở nên \emph{có ích} --- nó gói cả một mệnh đề thành một cái cũ để bắt đầu câu sau.
\end{cambay}

\section{Vì sao ``thêm từ nối'' chữa sai bệnh}

Câu hỏi mở đầu thứ ba nhắm vào lời khuyên phổ biến nhất về mạch văn. Từ nối --- \emph{however}, \emph{therefore}, \emph{moreover} --- có ích, nhưng chúng \emph{đặt tên} cho quan hệ giữa hai câu chứ không \emph{tạo} ra quan hệ đó.

Nếu thứ tự thông tin đã đúng, người đọc thường nhận ra quan hệ mà không cần từ nối\cite{halliday1976}; từ nối lúc đó chỉ xác nhận. Nếu thứ tự sai, thêm \emph{however} vào đầu một câu rời rạc không làm nó hết rời rạc --- nó chỉ tuyên bố có một sự tương phản mà người đọc không thấy.

Tệ hơn, từ nối đứng đầu câu \emph{chiếm mất} vị trí chủ đề. \emph{However, the residual increased} bắt đầu bằng bốn âm tiết không mang thông tin nội dung nào. Với những câu mà mối nối đã rõ, việc đó là lãng phí một vị trí quý.

Cách dùng đúng: dùng từ nối khi quan hệ \emph{không} suy ra được từ nội dung --- đặc biệt là tương phản và nhượng bộ, hai quan hệ mà người đọc không tự đoán ra. Bỏ chúng khi quan hệ đã hiển nhiên. Và nếu bạn thấy mình phải thêm từ nối vào mọi câu để đoạn đọc được, đó là dấu hiệu thứ tự thông tin sai chứ không phải thiếu từ nối.

\section{Mạch văn và bị động}

Câu hỏi mở đầu thứ tư nối lại với \ptr{C/11}. Ở đó ta nói bị động không sai; giờ ta có lý do chính xác nhất để dùng nó.

Giả sử câu trước vừa kết thúc bằng \emph{blurred frames}. Câu tiếp theo cần nói rằng hệ thống loại các khung đó. Bản chủ động --- \emph{The system discards such frames} --- bắt đầu bằng \emph{the system}, một thông tin cũ nhưng \emph{không phải cái cũ vừa nhắc}. Bản bị động --- \emph{Such frames are then discarded} --- bắt đầu đúng bằng mối nối.

Đây là điểm quan trọng: quyết định chủ động hay bị động không nên dựa trên một luật chung mà dựa trên \emph{cái gì cần đứng ở đầu câu để nối mạch}. Hai nguyên tắc của \ptr{C/11} và của chương này đôi khi kéo ngược nhau, và khi đó mạch văn thường thắng --- vì một câu hơi nặng mà nối đúng chỗ vẫn dễ đọc hơn một câu gọn mà làm người đọc mất mạch.

\section{Quy trình sửa một đoạn rời rạc}

Câu hỏi mở đầu thứ năm hỏi quy trình. Bốn bước, làm được một cách máy móc.

\emph{Bước một --- gạch dưới vài từ đầu của mỗi câu trong đoạn.} Chỉ phần trước động từ chính.

\emph{Bước hai --- đọc riêng chuỗi các phần đã gạch, bỏ qua phần còn lại.} Nếu chuỗi đó tự nó kể được một mạch --- cùng một chủ đề, hoặc một chuỗi móc xích --- thì đoạn ổn. Nếu nó là một danh sách các khái niệm rời rạc thì đó chính là bệnh.

\emph{Bước ba --- với mỗi câu bị lệch, hỏi: cái cũ trong câu này là gì?} Rồi viết lại để cái cũ đó lên đầu. Công cụ: đảo chủ động/bị động, tách câu, hoặc dùng danh từ gói lại ý trước.

\emph{Bước bốn --- kiểm vị trí nhấn.} Trong mỗi câu, thứ đáng nhớ nhất có đang ở cuối không? Nếu nó bị chôn giữa câu, hãy chuyển các cụm bổ ngữ lên trước.

\begin{capdoi}[Bước ba và bốn]
\truoc{Errors in the extrinsic calibration are the reason for the residual offset of about five millimetres that we observed.}
\sau{The residual offset we observed --- about five millimetres --- comes from errors in the extrinsic calibration.}
\vico{Câu trước vừa nói về offset, nên offset là cái cũ và phải lên đầu; nguyên nhân là cái mới nên về cuối, đúng vị trí nhấn.}
\truoc{The system, after discarding frames whose inlier count falls below the threshold, continues tracking.}
\sau{The system discards frames whose inlier count falls below the threshold, and then continues tracking.}
\vico{Bổ ngữ dài chen giữa chủ ngữ và động từ; người đọc chờ động từ quá lâu (\ptr{C/11}, phép bốn).}
\end{capdoi}

\section{Gốc rễ: từ trọng âm nói tới trật tự viết}

Ý ``cũ trước, mới sau'' không phải một quy ước văn phong do ai đó nghĩ ra. Nó bắt nguồn từ cách ngôn ngữ nói hoạt động, và các nhà ngôn ngữ học Praha đầu thế kỷ 20 là những người mô tả nó có hệ thống đầu tiên. Quan sát của họ: trong lời nói, câu có xu hướng đi từ phần ít mang thông tin mới nhất tới phần mang nhiều thông tin mới nhất, và trọng âm của câu rơi vào phần cuối đó. Thứ tự này không phải luật ngữ pháp mà là hệ quả của việc giao tiếp diễn ra theo thời gian: người nghe cần một điểm tựa trước khi nhận cái mới.

Truyền thống ngữ pháp chức năng ở Anh, gắn với Halliday\cite{halliday1976}, phát triển tiếp và tách bạch hai chiều: \emph{chủ đề} --- cái câu nói về, ở đầu --- và \emph{tiêu điểm thông tin} --- cái mới, ở cuối. Hai chiều này thường trùng nhau nhưng không luôn, và chính khoảng chênh giữa chúng giải thích vì sao đôi khi bạn phải dùng bị động hay đảo trật tự.

Đóng góp của Gopen và Swan năm 1990 không phải phát hiện ra hiện tượng mà là \emph{đổi câu hỏi}: thay vì hỏi ``câu này viết đúng chưa'', họ hỏi ``người đọc chờ gì ở vị trí nào, và văn bản có đáp đúng kỳ vọng đó không''. Cách đặt vấn đề này biến văn phong từ chuyện thẩm mỹ thành chuyện \emph{kỹ thuật}: có thể chẩn đoán, có thể sửa, và có thể dạy --- đúng như bốn bước ở mục 5.

Một hệ quả đáng nói: vì cơ chế bắt nguồn từ lời nói chứ không từ tiếng Anh nói riêng, nguyên tắc này áp dụng được cho tiếng Việt và nhiều ngôn ngữ khác. Cái thay đổi giữa các ngôn ngữ là \emph{công cụ} để đưa một thành phần lên đầu --- tiếng Anh dùng bị động và đảo ngữ, các ngôn ngữ khác có tiểu từ chủ đề hoặc trật tự tự do hơn --- chứ không phải bản thân nguyên tắc.

\begin{gocre}
\gr{Đầu thế kỷ 20 --- trường phái Praha}{Vì sao trật tự từ thay đổi dù nghĩa mệnh đề không đổi?}{Mô tả dòng thông tin trong câu: từ đã biết tới mới; trọng âm rơi ở cuối.}
\gr{Giữa thế kỷ 20 --- Halliday\cite{halliday1976}}{Câu làm mấy việc cùng lúc?}{Tách chủ đề (đầu câu) khỏi tiêu điểm thông tin (cuối câu) --- giải thích các trường hợp hai cái lệch nhau.}
\gr{1990 --- Gopen \& Swan\cite{gopenswan1990}}{Vì sao văn khoa học đúng ngữ pháp mà khó đọc?}{Đổi từ ``viết thế nào'' sang ``người đọc chờ gì ở đâu''; biến văn phong thành việc chẩn đoán được.}
\gr{Nguồn gốc lời nói}{Vì sao thứ tự này lại tự nhiên?}{Giao tiếp diễn ra theo thời gian; người nghe cần điểm tựa trước khi nhận cái mới --- nên nguyên tắc vượt ra ngoài tiếng Anh.}
\gr{Từ nối trong siêu diễn ngôn}{Từ nối làm gì?}{Đặt tên cho quan hệ, không tạo ra quan hệ --- nền cho mục 3 và cho \ptr{B/06} ở phía đọc.}
\end{gocre}

\section{Liên hệ ngang}

\begin{lienhe}
\lh{Thiết kế trình bày}{Slide và biểu đồ}{Cùng nguyên tắc vị trí: mắt người đi theo thứ tự dự đoán được, nên thông tin quan trọng đặt ở nơi mắt dừng lại. Kết luận của một biểu đồ nên nằm ở tiêu đề, không phải chú thích.}
\lh{Lập trình}{Thứ tự tham số và tên hàm}{API dễ dùng giữ thứ tự tham số nhất quán giữa các hàm --- cùng logic ``đừng phá kỳ vọng của người dùng theo vị trí''.}
\lh{Âm nhạc}{Câu nhạc và điểm rơi}{Trọng âm ở cuối câu nhạc; người nghe chờ nó. Vị trí nhấn trong văn có cùng cơ chế tâm lý.}
\lh{Tiếng Việt}{Cấu trúc ``thì, mà, là'' và trật tự chủ đề}{Tiếng Việt đánh dấu chủ đề bằng trật tự và tiểu từ chứ không bằng bị động; nguyên tắc giống, công cụ khác.}
\lh{Kể chuyện}{Mở đầu cảnh phim}{Cảnh mới thường mở bằng một khung hình nối với cảnh trước rồi mới giới thiệu cái mới --- móc xích ở cấp hình ảnh.}
\lh{Toán học}{Trình bày chứng minh}{Mỗi bước bắt đầu từ cái vừa thiết lập; một chứng minh khó theo thường là chứng minh nhảy cóc điểm tựa, không phải chứng minh sai.}
\end{lienhe}

\begin{traloi}
\tl{Ở \emph{thứ tự thông tin} trong từng câu. Đoạn dễ theo đặt cái người đọc đã biết ở đầu câu và cái mới ở cuối, nên mỗi câu mới đều có điểm tựa. Đoạn rời rạc mở mỗi câu bằng một khái niệm mới, khiến người đọc không thấy các câu liên quan gì với nhau --- dù từng câu đều đúng ngữ pháp và có đủ từ nối.}
\tl{Vì đầu câu là \emph{vị trí chủ đề}: người đọc dùng nó để tìm mối nối với những gì vừa đọc, không phải để nhận tin quan trọng. Chỗ dành cho cái đáng nhớ là cuối câu --- vị trí nhấn, nơi trọng âm rơi và trí nhớ giữ lại nhiều nhất. Lời khuyên ``đặt ý chính lên đầu'' đúng ở cấp đoạn và cấp bài (\ptr{C/13}) nhưng sai ở cấp câu.}
\tl{Vì từ nối \emph{đặt tên} cho quan hệ chứ không \emph{tạo} ra quan hệ. Nếu thứ tự thông tin đúng, người đọc thường thấy quan hệ mà không cần từ nối; nếu thứ tự sai, thêm \emph{however} chỉ tuyên bố một tương phản mà họ không thấy. Ngoài ra từ nối đứng đầu chiếm mất vị trí chủ đề. Dùng chúng khi quan hệ không suy ra được từ nội dung --- nhất là tương phản và nhượng bộ; bỏ khi đã hiển nhiên.}
\tl{Khi cái cần đứng đầu câu để nối với câu trước lại là \emph{đối tượng} của hành động. Nếu câu trước kết thúc bằng \emph{blurred frames}, thì \emph{Such frames are then discarded} nối tốt hơn \emph{The system discards such frames}, dù bản chủ động ngắn hơn. Khi nguyên tắc chủ thể--hành động của \ptr{C/11} và nguyên tắc mạch văn kéo ngược nhau, mạch văn thường thắng: câu hơi nặng mà nối đúng vẫn dễ hơn câu gọn mà làm mất mạch.}
\tl{Bốn bước. Gạch dưới phần trước động từ chính của mỗi câu. Đọc riêng chuỗi các phần đó: nếu nó tự kể được một mạch (cùng chủ đề, hoặc móc xích) thì đoạn ổn, nếu là danh sách khái niệm rời rạc thì đó là bệnh. Với mỗi câu lệch, xác định cái cũ và đưa nó lên đầu --- bằng đảo chủ động/bị động, tách câu, hoặc dùng danh từ gói lại ý trước. Cuối cùng kiểm xem thứ đáng nhớ nhất có nằm ở cuối câu không.}
\end{traloi}

\begin{dongian}
Có một điều về tiếng Anh mà ít ai nói cho người học, nhưng biết rồi thì sửa được rất nhiều bài viết: trong một câu, \emph{vị trí} của thông tin quan trọng gần bằng bản thân thông tin.

Đầu câu là chỗ người đọc tìm mối nối. Họ vừa đọc câu trước, và khi bước vào câu mới, mấy từ đầu tiên trả lời câu hỏi ``cái này liên quan gì tới cái vừa rồi?''. Nếu chỗ đó là một thứ hoàn toàn mới, họ hụt chân.

Cuối câu là chỗ nhấn. Đọc thầm một câu, bạn sẽ thấy giọng trong đầu lên trọng âm ở gần cuối. Cái nằm đó được nhớ lâu nhất.

Nên quy tắc rất gọn: \emph{cũ trước, mới sau}. Bắt đầu câu bằng thứ người đọc đã biết, kết thúc bằng thứ bạn muốn họ nhớ.

Thử xem nó hoạt động thế nào. Bốn câu này rời rạc: \emph{Motion blur degrades feature extraction. Fewer inliers are produced by the matching stage. A loss of tracking may follow. Aggressive rotation is the usual cause of blur.} Từng câu đúng cả, mà đọc như bốn mẩu tin không liên quan.

Bây giờ để cái cuối câu trước quay lại làm cái đầu câu sau: \emph{Motion blur degrades feature extraction. Poor features give the matching stage fewer inliers, and too few inliers can cause a loss of tracking.} Cùng nội dung, nhưng bây giờ nó chảy.

Còn một lời khuyên bạn nên bớt tin: ``thêm \emph{however}, \emph{moreover} cho mạch lạc''. Từ nối chỉ \emph{gọi tên} quan hệ, không tạo ra quan hệ. Nếu thứ tự đã đúng thì nhiều khi không cần chúng; nếu thứ tự sai thì thêm bao nhiêu cũng không cứu được.

Và một mẹo sửa nhỏ mà hiệu quả lớn: đừng viết \emph{This shows that\ldots} với \emph{this} trơ trọi. Cho nó một danh từ --- \emph{This drop shows\ldots}, \emph{This inconsistency shows\ldots} --- để người đọc khỏi phải quay lại đoán.
\motcau{Đầu câu để bắt tay, cuối câu để nhấn.}
\chuadongian{Khi nguyên tắc câu rõ và nguyên tắc mạch văn xung đột, chọn cái nào là một phán đoán tuỳ ngữ cảnh, không có luật.}
\end{dongian}

\begin{onbai}
\ar[3]{Hai vị trí, hai chức năng}{Đầu câu = vị trí chủ đề (nối với cái đã biết); cuối câu = vị trí nhấn (cái mới, cái đáng nhớ).}
\ar[3]{Quy tắc cũ trước mới sau}{Bắt đầu bằng thông tin người đọc đã có; kết thúc bằng thông tin bạn đang thêm.}
\ar[3]{Móc xích}{Phần mới ở cuối câu này thành phần cũ ở đầu câu sau --- dùng cho chuỗi nhân quả và quy trình.}
\ar[3]{Kiểu lặp lại}{Mọi câu bắt đầu bằng cùng một chủ đề --- dùng khi mô tả nhiều mặt của một đối tượng.}
\ar[2]{Đoạn rời rạc}{Không theo kiểu nào; đầu câu nhảy lung tung. Hoặc trộn hai kiểu giữa chừng không lý do.}
\ar[2]{Từ nối làm gì}{Đặt tên cho quan hệ, không tạo ra quan hệ; và chúng chiếm mất vị trí chủ đề.}
\ar[2]{Khi nào cần từ nối}{Khi quan hệ không suy ra được từ nội dung --- nhất là tương phản và nhượng bộ.}
\ar[2]{\emph{This} trơ trọi}{Luôn cho \emph{this} một danh từ gói lại ý trước: \emph{This drop}, \emph{This choice}.}
\ar[2]{Bị động vì mạch văn}{Dùng khi cái cần đứng đầu để nối lại là đối tượng của hành động.}
\ar[2]{Xung đột hai nguyên tắc}{\ptr{C/11} và \ptr{C/12} đôi khi kéo ngược; mạch văn thường thắng.}
\ar[2]{Bốn bước sửa đoạn}{Gạch đầu câu \(\rightarrow\) đọc riêng chuỗi đó \(\rightarrow\) đưa cái cũ lên đầu \(\rightarrow\) kiểm vị trí nhấn.}
\ar[1]{Nguồn gốc nguyên tắc}{Từ cách lời nói hoạt động --- người nghe cần điểm tựa trước khi nhận cái mới; nên nó vượt ra ngoài tiếng Anh.}
\end{onbai}

\begin{hoidap}
\qa{Quy tắc này có áp dụng cho tiếng Việt không?}{Có, vì nó bắt nguồn từ cách lời nói hoạt động chứ không riêng tiếng Anh. Cái khác là \emph{công cụ}: tiếng Anh đưa một thành phần lên đầu bằng bị động hoặc đảo ngữ, còn tiếng Việt dùng trật tự và các cấu trúc chủ đề kiểu ``Cái lỗi này thì\ldots''. Luyện trên tiếng Việt là cách rẻ để cảm nhận nguyên tắc, rồi mang sang tiếng Anh.}
\qa{Nếu mọi câu đều bắt đầu bằng cái cũ, văn có thành đều đều nhàm không?}{Rủi ro có thật nhưng nhỏ hơn bạn tưởng, vì ``cái cũ'' ở mỗi câu là một thứ khác nhau trong chuỗi móc xích. Chỗ dễ nhàm là kiểu lặp lại kéo dài: sáu câu liền mở bằng \emph{The system}. Cách phá đều: thỉnh thoảng mở bằng một cụm trạng ngữ ngắn đã biết (\emph{In the second experiment,\ldots}), hoặc gộp hai câu ngắn thành một câu có hai vế.}
\qa{Câu đầu tiên của một đoạn thì lấy đâu ra ``cái cũ''?}{Từ đoạn trước, hoặc từ tiêu đề phần. Câu mở đoạn có nhiệm vụ kép: nối với đoạn trước và tuyên bố chủ đề của đoạn mới. Nếu đoạn mới đổi hướng hẳn, hãy báo hiệu ra mặt (\emph{We now turn to\ldots}) --- đứt mạch có chủ ý thì người đọc chấp nhận, đứt mạch im lặng thì họ tưởng mình đọc sót.}
\qa{Tôi có nên viết theo quy tắc này ngay từ bản nháp đầu không?}{Không nên. Bản nháp đầu lo ghi ý ra; sắp lại thứ tự thông tin là việc của lượt sửa, và bốn bước ở mục 5 được thiết kế đúng cho lượt đó. Lý do thực dụng: khi viết nháp bạn chưa biết chắc câu sau sẽ nói gì, nên không thể tối ưu mối nối. \ptr{C/15} đặt bước này vào đúng chỗ trong quy trình.}
\qa{Có mâu thuẫn không giữa ``đặt cái mới ở cuối câu'' và lời khuyên ``đặt kết luận lên đầu bài''?}{Không, vì chúng ở hai cấp khác nhau. Ở cấp \emph{bài} và cấp \emph{đoạn}, đưa kết luận lên sớm là đúng --- người đọc cần biết bạn đang đi đâu (\ptr{C/13}). Ở cấp \emph{câu}, đầu câu phục vụ việc nối mạch. Một câu chủ đề đứng đầu đoạn vẫn nên tuân quy tắc bên trong nó: mở bằng mối nối, đóng bằng ý mới của đoạn.}
\qa{Làm sao tôi tự kiểm được mạch mà không cần người khác đọc?}{Dùng bước một và hai ở mục 5 --- chúng chính là bài kiểm tự làm. Một biến thể mạnh hơn: chép riêng phần đầu mỗi câu ra một danh sách và đọc to danh sách đó. Nếu nó nghe như một dàn ý có mạch thì đoạn ổn; nếu nghe như các mục lộn xộn thì bạn vừa tìm ra chỗ hỏng mà không cần đọc lại toàn bộ đoạn.}
\qa{Đọc bài của người khác, quy tắc này giúp gì?}{Giúp nhiều hơn bạn tưởng. Khi một đoạn khó theo, hãy thử bước hai: nếu chuỗi đầu câu lộn xộn thì khó khăn của bạn là do \emph{cách viết} chứ không do bạn thiếu nền. Điều này tiết kiệm rất nhiều tự nghi ngờ khi đọc tài liệu khó, và bổ sung cho \ptr{B/06}: ở đó bạn dựng lại lập luận, ở đây bạn chẩn đoán vì sao nó khó dựng.}
\end{hoidap}

\begin{baitap}
\bt[1]{Lấy một đoạn bạn viết, gạch dưới phần trước động từ chính của mỗi câu}{Bài viết của bạn}{Chép riêng chuỗi đó ra và đọc to.}
\bt[1]{Phân loại 5 đoạn trong một bài báo theo kiểu lặp lại hay móc xích}{Bài báo ngành}{Đoạn nào không thuộc kiểu nào thì đánh dấu và xem nó có khó theo không.}
\bt[2]{Sửa đoạn ở bài 1 theo bốn bước của mục 5}{Bài viết của bạn}{Không đổi nội dung, chỉ đổi thứ tự và cách nối.}
\bt[2]{Tìm 5 chỗ dùng \emph{this} trơ trọi trong bài của bạn và thêm danh từ cho mỗi chỗ}{Bài viết của bạn}{Chọn danh từ gói đúng ý câu trước --- đây là phần khó.}
\bt[2]{Lấy một đoạn có nhiều từ nối, xoá hết từ nối rồi đọc lại}{Bài báo ngành}{Chỗ nào mất từ nối mà vẫn hiểu thì từ nối đó thừa; chỗ nào hụt thì nó cần thiết.}
\bt[3]{Viết một đoạn 6 câu theo kiểu móc xích, rồi viết lại cùng nội dung theo kiểu lặp lại}{Tự chọn}{So hai bản: kiểu nào hợp với nội dung này hơn, và vì sao.}
\bt[3]{Đổi 5 câu chủ động thành bị động vì lý do mạch văn, và ghi lý do từng câu}{Bài viết của bạn}{Nếu không ghi được lý do mạch văn thì đừng đổi.}
\end{baitap}

\section*{Kết chương và đọc tiếp}

Hai chương vừa rồi cho hai nguyên tắc bù nhau: \ptr{C/11} lo bên trong câu, chương này lo chỗ nối giữa các câu. Cùng nhau chúng xử lý gần hết những gì người đọc gọi là ``văn khó theo''.

Nhưng vẫn còn một cấp nữa. Người đọc bài báo không chỉ theo mạch từng đoạn --- họ mang theo một bộ kỳ vọng về việc \emph{cả bài} phải được tổ chức thế nào: cái gì nằm ở phần mở đầu, tóm tắt phải trả lời những câu hỏi nào, một đoạn nên mở bằng gì. Những kỳ vọng đó là quy ước của cộng đồng, và biết chúng giúp bạn viết nhanh hơn nhiều. Chương \ptr{C/13} đi vào đó.
\ifSubfilesClassLoaded{\printbibliography[title={Nguồn}]}{}
\end{document}
